% ============================================================
% EconTeacher — Worksheet Template
% Compile from this directory so the relative \input path to
% econteacher.sty resolves correctly. Generated worksheets in
% resources/generated/worksheets/ use the same relative path.
%
% SECTION TYPES AVAILABLE:
%   Key Terms          — definitions with blank answer space
%   Group Discussion   — scenario boxes + dotfill response lines
%   Multiple Choice    — standard A–D MCQ questions
%   Short Response     — questions with answer space and marks
%   Quantitative       — calculation questions with working space
%   Diagram            — blank labelled axes for student drawing
%   Data Response      — extract box followed by questions
%   Essay              — extended response with evaluation tips
%
% ENVIRONMENTS (from econteacher.sty):
%   \begin{instructions}      — blue box for worksheet instructions
%   \begin{scenario}[TITLE]   — blue box for discussion scenarios
%   \begin{extract}[SOURCE]   — grey box for data response text
%   \begin{keypoint}          — yellow box for evaluation tips
% ============================================================

\documentclass[12pt,a4paper]{article}
\usepackage[margin=1in,headheight=15pt]{geometry}
\usepackage{enumitem}
\usepackage{amsmath}
\usepackage{fancyhdr}
\usepackage{booktabs}
\usepackage{array}
\usepackage{tikz}
\usepackage{pgfplots}
\pgfplotsset{compat=1.18}
\RequirePackage{tcolorbox}
\tcbuselibrary{skins}

% Key Point — yellow; for the single most important idea on a slide
\newtcolorbox{keypoint}{
    enhanced,
    colback=yellow!12,
    colframe=yellow!60!black,
    fonttitle=\bfseries,
    title=Key Point,
    arc=2pt,
    boxrule=0.8pt,
}

% Formula — blue; for named equations and their variable key
\newtcolorbox{formula}[1][Formula]{
    enhanced,
    colback=blue!5,
    colframe=blue!45!black,
    fonttitle=\bfseries,
    title=#1,
    arc=2pt,
    boxrule=0.8pt,
}

% Example Question — teal; for exam-style questions posed to students
\newtcolorbox{examplequestion}[1][Example Question]{
    enhanced,
    colback=teal!6,
    colframe=teal!55!black,
    fonttitle=\bfseries,
    title=#1,
    arc=2pt,
    boxrule=0.8pt,
}

% Example Solution — green; always follows an examplequestion block
\newtcolorbox{examplesolution}[1][Solution]{
    enhanced,
    colback=green!6,
    colframe=green!45!black,
    fonttitle=\bfseries,
    title=#1,
    arc=2pt,
    boxrule=0.8pt,
}

% ---- Worksheet environments --------------------------------

% Scenario — blue; for group discussion prompts on worksheets
% Usage: \begin{scenario}{Scenario 1: Title with commas, fine} ... \end{scenario}
\newtcolorbox{scenario}[1]{
    enhanced,
    colback=blue!5,
    colframe=blue!45!black,
    fonttitle=\bfseries,
    title={#1},
    arc=2pt,
    boxrule=0.8pt,
}

% Instructions — blue; for worksheet-level instruction blocks
\newtcolorbox{instructions}{
    enhanced,
    colback=blue!5,
    colframe=blue!45!black,
    fonttitle=\bfseries,
    title=Instructions,
    arc=2pt,
    boxrule=0.8pt,
}

% Extract — grey; for data response source material
% Usage: \begin{extract}{Source: Author, Year} ... \end{extract}
\newtcolorbox{extract}[1]{
    enhanced,
    colback=gray!8,
    colframe=gray!50,
    fonttitle=\bfseries,
    title={#1},
    arc=2pt,
    boxrule=0.8pt,
}


\pagestyle{fancy}
\fancyhf{}
\lhead{Dr Jac Thomas}
\rhead{Topic index: X.X.X.X}
\rfoot{Page \thepage}

\begin{document}

% ---- Title block -------------------------------------------
\begin{center}
    \Large\textbf{TOPIC TITLE}\\[5pt]
    \normalsize Topic index: X.X.X.X \quad|\quad AQA AS/A-level Economics\\[3pt]
    \textit{Total Marks: XX \quad Suggested Time: XX minutes}
\end{center}

\vspace{0.8em}

\begin{instructions}
    Answer all questions in the spaces provided. Show your reasoning clearly for calculation questions.
\end{instructions}

\vspace{1em}

% ============================================================
% KEY TERMS SECTION
% Use for vocabulary recall. 2 marks per definition is standard.
% Adjust \\ spacing to suit expected answer length.
% ============================================================
\section*{Key Terms (X marks)}

\textit{Define each of the following terms. Each definition is worth 2 marks.}

\vspace{0.5em}

\begin{enumerate}[leftmargin=*, label=\textbf{\arabic*.}]

    \item \textbf{TERM ONE} \\[3cm]

    \item \textbf{TERM TWO} \\[3cm]

    \item \textbf{TERM THREE} \\[3cm]

\end{enumerate}

\newpage

% ============================================================
% GROUP DISCUSSION SECTION
% Use for scenario-based group activities. Each scenario sits
% in a \begin{scenario}[TITLE] box. dotfill lines collect
% structured responses. State the discussion prompts once
% above all scenarios, not repeated per scenario.
% ============================================================
\section*{Group Discussion (XX minutes)}

\textit{In your groups, discuss each scenario. For each, decide:}
\begin{enumerate}[label=(\alph*), topsep=2pt]
    \item DISCUSSION PROMPT 1
    \item DISCUSSION PROMPT 2
    \item DISCUSSION PROMPT 3
\end{enumerate}

\vspace{0.8em}

\begin{scenario}{Scenario 1: TITLE}
    SCENARIO DESCRIPTION TEXT.
\end{scenario}

\vspace{0.5em}
\textbf{Your analysis:}\\
RESPONSE LABEL 1: \dotfill \\[0.4em]
RESPONSE LABEL 2: \dotfill \\[0.4em]
RESPONSE LABEL 3: \dotfill

\vspace{1em}

\begin{scenario}{Scenario 2: TITLE}
    SCENARIO DESCRIPTION TEXT.
\end{scenario}

\vspace{0.5em}
\textbf{Your analysis:}\\
RESPONSE LABEL 1: \dotfill \\[0.4em]
RESPONSE LABEL 2: \dotfill \\[0.4em]
RESPONSE LABEL 3: \dotfill

\newpage

% ============================================================
% MULTIPLE CHOICE SECTION
% Standard A–D options. Use (\Alph*) labels for options.
% 1 mark per question is typical; adjust section total accordingly.
% ============================================================
\section*{Multiple Choice (X marks)}

\textit{Circle the correct answer. Each question is worth 1 mark.}

\vspace{0.5em}

\begin{enumerate}[leftmargin=*, label=\textbf{\arabic*.}]

    \item QUESTION TEXT?
    \begin{enumerate}[label=(\Alph*), topsep=2pt]
        \item OPTION A
        \item OPTION B
        \item OPTION C
        \item OPTION D
    \end{enumerate}

    \vspace{0.4em}

    \item QUESTION TEXT?
    \begin{enumerate}[label=(\Alph*), topsep=2pt]
        \item OPTION A
        \item OPTION B
        \item OPTION C
        \item OPTION D
    \end{enumerate}

\end{enumerate}

\newpage

% ============================================================
% SHORT RESPONSE SECTION
% Questions requiring a few sentences to a paragraph. Use
% \\[Xcm] for shorter questions; \vspace{Xcm} for longer ones.
% Mark allocation appears inline at the end of the question text,
% before the answer space: QUESTION \hfill (X marks)\\[Xcm]
% ============================================================
\section*{Short Response (X marks)}

\textit{Answer all questions in the spaces provided.}

\vspace{0.5em}

\begin{enumerate}[leftmargin=*, label=\textbf{\arabic*.}]

    \item QUESTION TEXT. \hfill (X marks)\\[4cm]

    \vspace{0.5em}

    \item QUESTION TEXT. \hfill (X marks)\\[5cm]

\end{enumerate}

\newpage

% ============================================================
% QUANTITATIVE SKILLS SECTION
% For calculation questions. Always include a "Show your
% working" prompt. Sub-parts labelled (i), (ii) etc.
% ============================================================
\section*{Quantitative Skills (X marks)}

\textit{Show all calculations clearly. State the formula before substituting values.}

\vspace{0.5em}

\begin{enumerate}[leftmargin=*, label=\textbf{\arabic*.}]

    \item QUESTION TEXT. \hfill (X marks)
    \begin{enumerate}[label=(\roman*), topsep=4pt]
        \item SUB-QUESTION ONE
        \item SUB-QUESTION TWO
        \item SUB-QUESTION THREE
    \end{enumerate}

    \vspace{0.3em}
    \textit{Show your working:}
    \vspace{3cm}

    \vspace{0.5em}

    \item QUESTION TEXT. \hfill (X marks)\\
    \textit{Show your working:}
    \vspace{3cm}

\end{enumerate}

\newpage

% ============================================================
% DIAGRAM SECTION
% Provide blank labelled axes for the student to draw on.
% Set xtick=\empty and ytick=\empty so students add their own.
% State clearly what should be drawn and labelled.
% ============================================================
\section*{Diagram (X marks)}

\begin{enumerate}[leftmargin=*, label=\textbf{\arabic*.}]

    \item QUESTION TEXT — draw and label the diagram below. \hfill (X marks)

    \vspace{0.5em}
    \begin{center}
        \begin{tikzpicture}
            \begin{axis}[
                axis lines = middle,
                xlabel = {XLABEL},
                ylabel = {YLABEL},
                xmin=0, xmax=10, ymin=0, ymax=10,
                xtick=\empty, ytick=\empty,
                width=11cm, height=8cm,
                xlabel style={at={(axis description cs:1,0)}, anchor=north east, yshift=-10pt, font=\small},
                ylabel style={at={(axis description cs:0,1)}, anchor=south east, font=\small},
            ]
            % Blank — student draws here
            \end{axis}
        \end{tikzpicture}
    \end{center}

\end{enumerate}

\newpage

% ============================================================
% DATA RESPONSE SECTION
% Source text sits in an \begin{extract}[SOURCE] box.
% Questions follow as a numbered list beneath.
% Reference line/paragraph numbers in questions where relevant.
% ============================================================
\section*{Data Response (X marks)}

\begin{extract}{Source: PUBLICATION, YEAR}
    EXTRACT TEXT. This can span multiple paragraphs. Reference specific
    lines or data points in the questions below.

    \vspace{0.4em}
    SECOND PARAGRAPH OF EXTRACT.
\end{extract}

\vspace{0.8em}

\begin{enumerate}[leftmargin=*, label=\textbf{\arabic*.}]

    \item Using the data, QUESTION TEXT. \hfill (X marks)\\[4cm]

    \item QUESTION TEXT. \hfill (X marks)\\[5cm]

\end{enumerate}

\newpage

% ============================================================
% ESSAY SECTION
% For extended responses (8+ marks). The \begin{keypoint} box
% gives evaluation prompts. Use \vspace to allocate writing
% space proportional to marks.
% ============================================================
\section*{Essay (X marks)}

\textit{Write a well-structured response. Define key terms, develop chains of reasoning, and evaluate your arguments.}

\vspace{0.5em}

\begin{enumerate}[leftmargin=*, label=\textbf{\arabic*.}]

    \item ESSAY QUESTION. \hfill (X marks)\\[0.8em]
    \begin{keypoint}
        \textbf{Evaluation tips:} TIP 1. Consider TIP 2. Think about TIP 3.
    \end{keypoint}
    \vspace{12cm}

\end{enumerate}

\vfill
\begin{center}
    \textbf{[Total: XX marks]}
\end{center}

\end{document}
